\documentclass[a4paper,12pt]{article}
\usepackage[utf8]{inputenc}
\usepackage[T1]{fontenc}
\usepackage[spanish]{babel}
\usepackage{amsmath, amsfonts, amssymb}
\usepackage{graphicx}
\usepackage{enumitem}
\usepackage{geometry}
\geometry{left=2cm,right=2cm,top=2cm,bottom=2cm}

\begin{document}

\begin{center}
    \Large \textbf{Evaluación recuperatorio de Electromagnetismo e Interacción Electromagnética} \\
    \large Escuela Merín, 6to Año \\
    \today
\end{center}
\vspace{1cm}

\section*{Instrucciones}
Resuelva cada uno de los problemas propuestos mostrando todos los pasos necesarios. Utilice las fórmulas y constantes oficiales según lo visto en clase. Se permite el uso de calculadora científica no celular.



\section*{Tema 1: Electromagnetismo}

\begin{enumerate}[label=\textbf{1.\arabic*}]
    \item \textbf{Electroimán para levantamiento:}  
    Un electroimán con núcleo de chapa de silicio posee $N=150$ espiras y funciona con una corriente de $I=0.4\,\mathrm{A}$. Se desea levantar una lata de conservas de hojalata (sin recubrimiento apreciable). Con base en la fuerza magnetomotriz y las características del material, determine si el electroimán es capaz de levantar la lata. (Utilice datos experimentales y relacione la FEM con la fuerza de atracción).

    \item \textbf{Fuerza de atracción:}  
    Calcular la fuerza de atracción de un electroimán de armadura de hierro si la inducción magnética en el núcleo es de $B=1.3\,\mathrm{T}$ y el área de contacto entre el núcleo y el hierro móvil es de $42\,\mathrm{cm}^2$. (Recuerde convertir $42\,\mathrm{cm}^2$ a $\mathrm{m}^2$).

    \item \textbf{Corriente requerida:}  
    Se requiere que un electroimán genere una fuerza de atracción de $2\,\mathrm{kN}$. Si el núcleo está fabricado de hierro forjado y se conoce que posee $N=2000$ espiras, determine la intensidad de corriente necesaria para alcanzar dicha fuerza, empleando las fórmulas oficiales para fuerza magnética y magnetomotriz.

    \item \textbf{Permeabilidad absoluta:}  
    Calcular la permeabilidad absoluta $\mu$ de una lámina de hierro que, al ser sometida a una intensidad de campo de $H=400\,\mathrm{A/m}$, adquiere una inducción magnética de $B=1.4\,\mathrm{T}$. Utilice la relación 
    \[
    \mu = \frac{B}{H}.
    \]

    \item \textbf{Inducción en chapa de silicio:}  
    Una chapa de silicio es sometida a un campo magnético de $H=9000\,\mathrm{A/m}$. Utilizando la curva característica oficial para este material (según datos experimentales), determine el valor de la inducción magnética $B$ en Teslas.
\end{enumerate}

\newpage
\begin{center}
    \Large \textbf{Evaluación de Electromagnetismo e Interacción Electromagnética} \\
    \large Escuela Merín, 6to Año \\
    \today
\end{center}
\vspace{1cm}
\section*{Instrucciones}
Resuelva cada uno de los problemas propuestos mostrando todos los pasos necesarios. Utilice las fórmulas y constantes oficiales según lo visto en clase. Se permite el uso de calculadora científica, no celular.
\section*{Tema 2: Interacción Electromagnética}

\begin{enumerate}[label=\textbf{2.\arabic*}]
    \item \textbf{FEM inducida por movimiento:}  
    Un conductor lineal de longitud $l=1\,\mathrm{m}$ se desplaza perpendicularmente a un campo magnético de inducción $B=1.5\,\mathrm{T}$ con una velocidad de $v=10\,\mathrm{m/s}$. Calcular la fuerza electromotriz inducida utilizando la fórmula 
    \[
    \varepsilon = B \, l \, v.
    \]

    \item \textbf{Velocidad para FEM requerida:}  
    ¿A qué velocidad debe moverse un conductor de longitud $l=0.5\,\mathrm{m}$ en un campo de inducción $B=1\,\mathrm{T}$ para que se induzca una fuerza electromotriz de $\varepsilon = 5\,\mathrm{V}$?

    \item \textbf{Autoinducción en bobina:}  
    Una bobina tiene un coeficiente de autoinducción de $L=50\,\mathrm{mH}$. Si se aplica una corriente que varía uniformemente de $0\,\mathrm{A}$ hasta $I=20\,\mathrm{A}$ en un intervalo de tiempo de $\Delta t=0.1\,\mathrm{s}$, determine la fuerza electromotriz de autoinducción que se desarrolla, usando la fórmula 
    \[
    \varepsilon_{\text{auto}} = -L \frac{\Delta I}{\Delta t}.
    \]

    \item \textbf{Determinación de número de espiras:}  
    Una bobina experimenta una variación de flujo magnético de $\Delta \Phi = 0.01\,\mathrm{Wb}$ cuando circula una corriente de $10\,\mathrm{A}$, y su coeficiente de autoinducción es $L=0.5\,\mathrm{H}$. Utilice las relaciones entre flujo, autoinducción y magnetomotriz para determinar el número de espiras $N$ de la bobina.

    \item \textbf{Diseño de generador:}  
    Diseñe un generador que debe producir una fuerza electromotriz de $\varepsilon = 10\,\mathrm{V}$. Suponga que un conductor se desplaza a una velocidad de $v=3.14\,\mathrm{m/s}$ en un campo magnético uniforme. Propuesta: determine, justificadamente, un número de espiras y la configuración geométrica (por ejemplo, en forma de solenoide) que permita alcanzar la FEM requerida, utilizando los conceptos de inducción electromagnética.
\end{enumerate}

\end{document}
